\section{Reflection normal closures and the remaining finite-group source gate}
\label{sec:reflection-reduction}

We work with a faithful absolutely irreducible representation over an
imaginary quadratic field \(K\), whose character field is exactly \(K\).
This is the split group-algebra block \(M_m(K)\) relevant to a generic
Weil factor with endomorphism algebra \(K\). The following argument is
not a classification of representations without reflections.

\subsection*{D-FIELD-01. The reflection components over the seven fields}

Let
\[
 \mathcal D=\{2,7,11,19,43,67,163\},\qquad K_d=\mathbb Q(\sqrt{-d}).
\]
An irreducible finite complex reflection group with a model over \(K_d\)
is a rational Weyl reflection representation, with the following sole
additional possibilities:
\[
  d=2:\ G_{12}\text{ in rank two};\qquad
  d=7:\ G_{24}\text{ in rank three}.
\]

Here we use the published reflection classification, rather than claim
to reprove it. Precise inputs are Cohen's rank-two classification and
table, \cite[\S\S3.4--3.6, pp.~394--395]{FIELD:Cohen}, and his
higher-rank classification, \cite[Theorem 5.12, p.~435]{FIELD:Cohen}.
The exceptional character fields are listed in
\cite[Tables 8.1--8.2, pp.~674--675]{FIELD:Malle}.
The table omits \(G_4,G_{25},G_{28},G_{35}\), as explicitly explained
on its following page. The first two have reflections of order three
(Cohen's \S3.6 and \S4.16 tables, pp.~395 and 412), so cannot have
models over these \(K_d\); the last two are the rational Weyl groups
\(F_4,E_6\). Inspection of the other listed fields leaves precisely
the two nonreal cases displayed above. Noncrystallographic real types
\(H_3,H_4\) have character field \(\mathbb Q(\sqrt5)\), and are excluded.

For completeness, the infinite imprimitive series needs no field table.
In \(G(a,b,k)\), \(k\ge2\), the diagonal element with entries
\(\zeta_a,\zeta_a^{-1},1,\ldots,1\) shows
\(2\cos(2\pi/a)\in K_d\cap\mathbb R=\mathbb Q\). Thus
\(a\in\{1,2,3,4,6\}\): the sum is a rational algebraic integer in
\([-2,2]\), and its five possible integer values give that list.
If \(k\ge3\), multiplying this diagonal element by the transposition
of the second and third coordinates gives trace
\(\zeta_a+k-3\), so \(\zeta_a\in K_d\) and \(a\le2\).
For \(k=2,b=a\), the group is the dihedral Weyl representation of
type \(A_2,B_2,G_2\), or a reducible lower case. For \(k=2,b<a\),
the element \(\operatorname{diag}(\zeta_a^b,1)\) forces
\(a/b\le2\). The remaining possibilities \((a,b)=(4,2),(6,3)\)
have traces \(2i\) and \(\zeta_6+\zeta_6^2=\sqrt{-3}\), respectively,
so are excluded. The cases \(a\le2\) are rational Weyl cases.
Rank one similarly has only the reflection \(-1\). The real standard
representations of symmetric groups are also rational Weyl cases.

\subsection*{D-FIELD-02. Rational normalizers and induced blocks}

\begin{lemma}\label{lem:weyl-normalizer}
Let \(W\) be an irreducible rational Weyl reflection group on
\(K_d^k\), with its standard rational positive definite form \(B\).
Let \(J\) be a finite subgroup of \(U(B,K_d)\) normalizing \(W\).
If \(d\in\{7,11,19,43,67,163\}\), every element of \(J\) has rational
matrix entries. If \(d=2\) and \(k\ge3\), either the same conclusion
holds or \(W\) is of type \(F_4\) and, up to the fixed rational
root coordinates, \(J\subset G_F\), the 2304-element group of D-GF4-01.
If \(J\) contains \(W\) and has a nonrational element in this last
case, then \(J=G_F\).
\end{lemma}

\begin{proof}
For \(g\in J\), conjugation by \(g\) is an automorphism of \(W\).
The equations
\[
                  A w=(g w g^{-1})A\qquad(w\in W)
\]
have rational coefficients. Their solution space over \(K_d\) is
one-dimensional, by absolute irreducibility. Hence it has a nonzero
rational solution \(A\), which is invertible, and \(g=cA\) for
some \(c\in K_d^*\). A finite order \(N\) of \(g\) gives
\(A^N=c^{-N}I\), so \(c^N\in\mathbb Q\). Consequently
\((c/\bar c)^N=1\). These fields have only the roots of unity
\(\pm1\), so either \(c\in\mathbb Q\) or
\(c\in\sqrt{-d}\,\mathbb Q\).

Unitarity implies
\[
                    A^tBA=\nu B,\qquad N(c)\nu=1,
                    \qquad \nu\in\mathbb Q_{>0}.
\]
The rational matrix \(A\) permutes the reflection root lines. For
one rational root \(\alpha\), write \(A\alpha=q\beta\), with
\(q\in\mathbb Q^*\) and \(\beta\) another standard root. Then
\[
             \nu=q^2\frac{B(\beta,\beta)}{B(\alpha,\alpha)}.
\]
An irreducible crystallographic root system has one root length or
two root lengths whose squared ratio is two or three. Thus the positive
rational square class of \(\nu\) is among \(1,2,3\).
If \(c\) is purely imaginary, its norm has square class \(d\).
The identity \(N(c)\nu=1\) therefore excludes that case for
\(d=7,11,19,43,67,163\), proving rationality of \(g\).

For \(d=2\), a nonrational \(g\) requires a root-length interchange
with squared length ratio two. In rank at least three the only
possibilities are \(B_k\) and \(F_4\). The two root-line conjugacy
classes in \(B_k\) have sizes \(k\) and \(k(k-1)\). A normalizer
cannot interchange them when \(k\ge3\), because their sizes differ.
Hence only \(F_4\) remains.

Use the standard \(F_4\) roots in \(\mathbb Q^4\) and the matrix
\(\Phi\) of D-GF4-01. The orthogonal matrix
\(\Phi/\sqrt2\) interchanges the two root-line classes, as was
checked directly on all 48 roots. An orthogonal normalizer preserving
the two classes preserves the actual root system. Such an automorphism
belongs to \(W(F_4)\): it sends a Weyl chamber to a Weyl chamber,
so after composition with \(W(F_4)\) it fixes a chamber and induces
an automorphism of its Cartan diagram; the oriented \(F_4\) diagram
has no nonidentity automorphism. It follows that a rational similarity
normalizing \(W(F_4)\) is a rational scalar times either a Weyl
element or a Weyl element times \(\Phi\).

In the first case a finite unitary element over \(K_2\) is just a
Weyl element, since its scalar is a root of unity in \(K_2\).
In the second case write it as \(\lambda w\Phi\). Unitarity gives
\(2N(\lambda)=1\), and the preceding finite-order argument gives
that \(\lambda\) is rational or purely imaginary. The rational
case would require a rational square equal to \(1/2\); hence
\(\lambda=\pm\sqrt{-2}/2\). Thus every such element lies in
\(W(F_4)z\), and \(J\subset W(F_4)\sqcup W(F_4)z=G_F\).
If \(J\) contains \(W(F_4)\) and one element of the second coset,
it contains both whole cosets and equals \(G_F\).
\end{proof}

Now let \(G\subset\operatorname{GL}_m(K_d)\) be finite and absolutely
irreducible, with exact character field \(K_d\), and suppose it contains
a reflection. Let \(R\) be the normal subgroup generated by all its
reflections. The fixed space \(V^R\) is \(G\)-stable, hence zero.
The nonorthogonality graph of the reflection root lines gives an
orthogonal decomposition
\[
                  V=V_1\oplus\cdots\oplus V_s
\]
into the irreducible reflection components of \(R\). Each root line
is defined over \(K_d\), being the image of a rank-one matrix
\(r-1\), so each \(V_i\) is defined over \(K_d\). The group \(G\)
permutes these components transitively, by irreducibility. If \(H\)
stabilizes \(V_1\), and \(\tau\) is its representation on \(V_1\),
then
\[
                         \rho=\operatorname{Ind}_H^G\tau.
\]
The image of \(H\) is finite, contains the irreducible reflection
group \(R_1\), and normalizes it. The character of \(\tau\) cannot
be rational-valued, since otherwise its induced character would be
rational-valued.

If \(R_1\) is a rational Weyl representation, a \(K_d\)-change of
basis puts it into its rational standard model. Such a change of basis
exists because the intertwining equations are defined over \(K_d\)
and the representations are absolutely irreducible. Its invariant
Hermitian form is a positive rational scalar times the standard rational
form. Lemma~\ref{lem:weyl-normalizer} therefore applies to the image of
\(H\).

We obtain the following conclusions.
\begin{enumerate}
\item For \(d=11,19,43,67,163\), the existence of a reflection
is impossible. Every component is a rational Weyl component by
D-FIELD-01, and the lemma would make \(\tau\), hence \(\rho\),
rational-valued.
\item For \(d=7\), the component must be \(G_{24}\) of rank three.
Its finite unitary normalizing overgroups are contained in scalar multiples of
\(G_{24}\), by Cohen's corrected Proposition 5.14
\cite[p.~435]{FIELD:Cohen}\cite[p.~613]{FIELD:CohenErr}.
The unique exception in that corrected proposition is \(G_{25}\),
which is irrelevant here. A scalar multiplier belonging to \(K_7\)
and giving a finite overgroup is \(\pm1\); both scalars already
belong to \(G_{24}\). Thus the image of \(H\) is exactly \(G_{24}\).
\item For \(d=2\), either \(\dim_{K_2}V_1=2\), or the image of
\(H\) is exactly \(G_F\) in rank four. This uses D-FIELD-01 for
the nonreal reflection case \(G_{12}\), and the lemma for every
real Weyl case of rank at least three. Rank one cannot have a
nonrational finite character over \(K_2\).
\end{enumerate}

\subsection*{D-FIELD-03. Consequences for dominant sixfold curve suppliers}

Suppose a family of connected finite \(G\)-covers produces a Weil
sixfold as the multiplicity factor for \(\rho\). For the block
stabilizer above, put \(Y=C/H\). Shapiro's isomorphism identifies
the \(\rho\)-multiplicity Hodge structure in \(H^1(C)\) with the
\(\tau\)-multiplicity Hodge structure for the \(H\)-cover \(C\to Y\).
One can see this directly by assigning a vector in each induced coset
summand to the corresponding lift of a singular simplex; the boundary
maps and Hodge de Rham complexes agree. The cup pairings agree with
the sum over the sheets, using the induced invariant Hermitian form.
The originally chosen irreducible invariant form differs from it by a
positive rational scalar. Thus this is an identification of polarized
rational variations up to a fixed rational scalar, and preserves the
signed determinant class in even dimension. Quotienting the kernel
of \(\tau\) gives a connected cover with the indicated faithful
image group. In particular the period map factors through the appropriate
Hurwitz space of these covers of \(Y\).

For a rank-\(k\) local system on a base of genus \(\gamma\), with
moving-dimension sum \(D\), a sixfold multiplicity space satisfies
\[
                       6=(2\gamma-2)k+D.
\]
There are at most \(D\) genuine branch points. In rank two this
gives \(D=10-4\gamma\). The possible Hurwitz dimensions are at
most seven for \(\gamma=0\), six for \(\gamma=1\), and five for
\(\gamma=2\); higher genera are impossible. Consequently an induced
rank-two source cannot have period rank nine.

For \(K_7\), D-FIELD-02 reduces every source containing a reflection
to the natural \(G_{24}\) representation on the intermediate base
\(Y\). D-DET-02 proves that any rank-nine image of this family has
split determinant. For \(K_2\), rank two is excluded by the preceding
dimension bound; the remaining possibility is the natural \(G_F\)
representation, and D-DET-03 again forces split determinant.

\textbf{Universal reflection-source conclusion.}
Over \(\mathbb Q(\sqrt{-2})\) or \(\mathbb Q(\sqrt{-7})\), every
dominant sixfold finite-group curve supplier whose faithful absolutely
irreducible representation has exact character field \(K\) and contains
a reflection lies in the split determinant component. This statement
allows arbitrary representation dimension, finite group, initial branch
orders, and base genus within that class.

For the five fields \(K_d\), \(d=11,19,43,67,163\), such a
representation contains no reflection at all. A possible dominant
source over those fields, or a nonsplit dominant source over \(K_2,K_7\),
must therefore use a representation without reflections. If its degree
is \(m\), each nonidentity inertia has moving dimension at least two.
A genus-one base has \(D=6\) and at most three branch points, hence
cannot give rank nine. On a base of genus at least two, the dimension
equation and \(m\ge2\) leave only \(\gamma=2,m\le3\), whose
Hurwitz dimension is less than nine. Thus the base must have genus zero,
and
\[
                 12\le r\le\frac{2m+6}{2}=m+3,
                 \qquad m\ge9.
\]
More generally, if every nonidentity element has moving rank at least
\(q\), the necessary inequality is \(m\ge6q-3\).

These are necessary conditions only. They neither produce a rank-nine
period map nor exclude all representations without reflections. The
remaining finite-group gate is precisely a genus-zero representation
of degree at least nine without reflections, with an actual rank-nine
period computation and a nonnorm value in D-DET-01 for a nonsplit target.
